% section: パターン別の思考フロー

\section{パターン別の思考フロー}

第2章では,漸化式を型ごとに解いた.ここでは,各型の公式を繰り返す
のではなく,問題を見たときにどの方向へ進むかをまとめる.

\subsection{基本パターン}

\begin{description}
  \item[等差型]
    $a_{n+1}-a_n$ を計算する.一定なら等差型,
    $n$ の式なら階差を足し戻す.
  \item[等比型]
    $a_{n+1}/a_n$ を計算する.一定なら等比型,
    $n$ の式なら比を掛け合わせる.
  \item[一次式型]
    $a_{n+1}=pa_n+q$ なら不動点を引く.
  \item[階差型]
    $a_{n+1}-a_n=f(n)$ なら,
    \[
      a_n=a_1+\sum_{k=1}^{n-1}f(k)
    \]
    へ直す.
  \item[階比型]
    $a_{n+1}/a_n=g(n)$ なら,
    \[
      a_n=a_1\prod_{k=1}^{n-1}g(k)
    \]
    へ直す.
\end{description}

\subsection{応用パターン}

\begin{description}
  \item[補助数列型]
    $f(n)$ が残る一次漸化式では, $f(n)$ と同じ種類の特解を探して
    引く.
  \item[隣接3項型]
    過去2項が線形に現れたら, $a_n=\lambda^n$ を代入して
    特性方程式を作る.
  \item[連立型]
    複数の数列が線形に更新されるなら,係数を組み合わせて
    一つの数列が等比型になる方向を探す.
  \item[分数型]
    逆数,不動点,不動点からの比を順に試す.
  \item[和を含む型]
    $S_n=a_1+\cdots+a_n$ と置き,元の数列と $S_n$ の関係を作る.
  \item[母関数型]
    畳み込みや隣接項の和が繰り返されるなら,係数列を一つの関数へ
    まとめ,和を積へ翻訳する.
\end{description}

\subsection{特殊関数型}

多項式の反復が現れたら,展開を続ける前に公式を探す.

\begin{description}
  \item[三角関数]
    $2x^2-1$ なら $x=\cos\theta$ を試す.
    $2x/(1-x^2)$ なら $x=\tan\theta$ を試す.
  \item[基本対称式]
    $x^2-2$ なら $x=u+u^{-1}$ を試す.
    \[
      (u+u^{-1})^2-2=u^2+u^{-2}.
    \]
  \item[ニュートン法]
    接線の交点を次の項にしている式なら,一般項ではなく誤差の減少を
    調べる.
  \item[初項依存型]
    初項に特別な値が指定されているなら,その値が恒等式や固定点を
    作っていないかを確認する.
\end{description}

ここでは「特殊な公式を知っているか」だけが勝負ではない.
置換後も同じ表現が保たれるかを計算し,保たれなければ別の方法へ戻る.

\subsection{解法を切り替える目安}

次のどれかが起こったら,今の方法を一度止める.

\begin{enumerate}
  \item 差を取ったのに,元の式より複雑になった.
  \item 比を取ったが,分母が0になる可能性を処理できない.
  \item 置換後に新しい種類の項が増え続ける.
  \item 初項を説明できない候補になっている.
\end{enumerate}

この手順を身につけると,個別のパターンを暗記するだけでなく,
見慣れない漸化式にも最初の一手を選べるようになる.
